\subsection{The Gamma--Fourier--Tate boundary construction}
\label{ss:gammafouriertate}

The local analysis above proves the primitive inequality on a nontrivial
initial interval and identifies the obstruction to propagation as a shorted
old/new coupling.  We now construct the arithmetic and archimedean boundary
data entering that coupling.  The construction is independent of the zeros
of $\xi$.  It fixes the amplitude, the moving Tate constraints and the
positive Gamma storage before the remaining estimate is formulated.

\subsubsection{The finite logarithmic boundary port}

Let $\mathcal H_N$ be the Hilbert space with orthonormal basis
$\{\delta_d:1\leq d\leq N\}$.  For $n\leq N$, let
$S_n\delta_d=\delta_{nd}$ when $nd\leq N$ and let it be zero otherwise.
For $t>0$ put
\begin{equation}\label{eq:newdZetaConnection}
 \mathscr Z_t=\sum_{n\leq N}n^{-t/2}S_n,
 \qquad
 \mathscr C_t=\sum_{n\leq N}\Lambda(n)n^{-t/2}S_n,
 \qquad Q\delta_d=(\log d)\delta_d.
\end{equation}
All the sums in this subsection belong to the finite divisibility algebra.
In particular, no convergence or domain assertion is implicit.

\begin{lemma}[Finite logarithmic connection]
\label{lem:newdFiniteConnection}
One has
\begin{align}
 [Q,\mathscr Z_t]&=\mathscr Z_t\mathscr C_t
                  =\mathscr C_t\mathscr Z_t,
                  \label{eq:newdCommutator}\\
 \mathscr Z_t'&=-\tfrac12\mathscr Z_t\mathscr C_t,
                  \label{eq:newdZetaDerivative}\\
 (\mathscr Z_t^{-1})'&=\tfrac12\mathscr C_t\mathscr Z_t^{-1}.
                  \label{eq:newdMobiusDerivative}
\end{align}
\end{lemma}

\begin{proof}
Since $[Q,S_n]=(\log n)S_n$, the coefficient of $S_k$ in the
commutator is $(\log k)k^{-t/2}$.  The coefficient of $S_k$ in
$\mathscr Z_t\mathscr C_t$ is
\[
 k^{-t/2}\sum_{d\mid k}\Lambda(d)=k^{-t/2}\log k.
\]
The algebra is commutative, which proves \eqref{eq:newdCommutator}.
Termwise differentiation proves \eqref{eq:newdZetaDerivative};
differentiating $\mathscr Z_t^{-1}\mathscr Z_t=I$ then proves
\eqref{eq:newdMobiusDerivative}.
\end{proof}

We record the operator identity for which this connection is needed.  Let
$G(t)$ be the positive Gram family at a threshold, let
$v(t)=(-h(t),I)^t$ be its harmonic residual, and let $P_0$ be the old
coordinate projection.  Write $M_0$ for the old block of $G(t_0)$.  The
normal equation is $P_0G(t)v(t)=0$.  Hence differentiation gives
\begin{equation}\label{eq:newdMovingNormal}
 P_0G'(t_0)v(t_0)=M_0h'(t_0).
\end{equation}
Suppose, as in the localized explicit-formula decomposition, that
\[
 G'=R^*\mathscr LR-\tfrac12(QG+GQ),
 \qquad Gv=S\delta_1,
\]
where $S>0$ is the scalar normalization.  With
\[
 W_{\rm H}=\mathscr L^{1/2}Rv,
 \qquad H_0=\mathscr L^{1/2}RP_0^*,
\]
equation \eqref{eq:newdMovingNormal} becomes
\begin{equation}\label{eq:newdFiniteExcess}
 H_0^*W_{\rm H}=M_0h'+B_Q,
 \qquad
 B_Q=\frac S2P_0
 \bigl[-\mathscr C\delta_1+\mathscr Z\mathscr C^*m\bigr],
 \qquad m=\mathscr Z^{-1}\delta_1,
\end{equation}
where the parameter is now fixed and omitted from the notation.

Define the doubled logarithmic boundary port by
\begin{align}
 W_\infty
 &=\sqrt{\frac S2}
   \binom{\mathscr C\delta_1}{\mathscr C^*m},
   \label{eq:newdWinfty}\\
 H_{\infty,0}x
 &=\sqrt{\frac S2}
   \binom{P_0^*x}{-\mathscr Z^*P_0^*x}.
   \label{eq:newdHinfty}
\end{align}

\begin{theorem}[Exact finite Gamma--Tate amplitude]
\label{thm:newdAmplitude}
The port \eqref{eq:newdWinfty}--\eqref{eq:newdHinfty} is defined entirely
by the finite Euler transport and satisfies
\begin{equation}\label{eq:newdAmplitude}
 H_{\infty,0}^*W_\infty=-B_Q.
\end{equation}
Consequently
\begin{equation}\label{eq:newdCompletedDivergence}
 (H_0^*\oplus H_{\infty,0}^*)
 (W_{\rm H}\oplus W_\infty)=M_0h'.
\end{equation}
\end{theorem}

\begin{proof}
Taking the adjoint of \eqref{eq:newdHinfty} and applying it to
\eqref{eq:newdWinfty} gives
\[
 H_{\infty,0}^*W_\infty
 =\frac S2P_0
   (\mathscr C\delta_1-\mathscr Z\mathscr C^*m)
 =-B_Q
\]
by \eqref{eq:newdFiniteExcess}.  Adding this equality to
\eqref{eq:newdFiniteExcess} proves
\eqref{eq:newdCompletedDivergence}.
\end{proof}

The theorem fixes both orientations and the normalization of the boundary
amplitude.  It is an amplitude identity, not an energy inequality: its
finite-dimensional port cannot be identified with the complete
infinite-rank Hodge factor.  The continuum realization below uses it only
for the arithmetic cross-incidence that it computes.

\subsubsection{The moving Tate terminal}

Let $H_0,E$ be Hilbert spaces, let $B_0:H_0\to\C^r$ be bounded and onto,
and let $B_E:E\to\C^r$ be bounded.  Put
\begin{equation}\label{eq:newdMovingData}
 K_0=\ker B_0,
 \quad G_0=B_0B_0^*,
 \quad C=B_0^*G_0^{-1}B_E,
 \quad M=I_E+C^*C=I_E+B_E^*G_0^{-1}B_E.
\end{equation}
The enlarged moment map is $B(h,e)=B_0h+B_Ee$.

\begin{theorem}[Conservative update of a moving constraint]
\label{thm:newdMovingConstraint}
There is an orthogonal decomposition
\begin{equation}\label{eq:newdMovingDecomposition}
 \ker B=(K_0\oplus0)\ \widehat\oplus\
          \{(-Ce,e):e\in E\}.
\end{equation}
Moreover
\begin{equation}\label{eq:newdMovingIsometry}
 W:E\longrightarrow H_0\oplus E,
 \qquad
 Wx=\binom{-CM^{-1/2}x}{M^{-1/2}x}
\end{equation}
is an isometry onto the second summand.  In particular,
\begin{equation}\label{eq:newdMovingConservation}
 \|M^{-1/2}x\|^2+\|CM^{-1/2}x\|^2=\|x\|^2.
\end{equation}
\end{theorem}

\begin{proof}
The surjectivity of $B_0$ gives
$H_0=K_0\widehat\oplus\operatorname{Ran} B_0^*$.  If $h=k+r$ is the corresponding
decomposition, then $B(h,e)=0$ is equivalent to $B_0r=-B_Ee$ and hence to
$r=-B_0^*G_0^{-1}B_Ee=-Ce$.  This proves
\eqref{eq:newdMovingDecomposition}.  Finally,
\[
 \|Wx\|^2
 =\langle M^{-1/2}x,(I+C^*C)M^{-1/2}x\rangle=\|x\|^2,
\]
which proves both remaining assertions.
\end{proof}

For $I_T=(-T,T)$ take the two Tate moments
\begin{equation}\label{eq:newdTateMomentMap}
 B_Tf=
 \binom{\langle f,e^{t/2}\rangle}{\langle f,e^{-t/2}\rangle}.
\end{equation}
Their Gram is
\begin{equation}\label{eq:newdTateGram}
 B_TB_T^*=2
 \begin{pmatrix}\sinh T&T\\T&\sinh T\end{pmatrix},
\end{equation}
with eigenvalues $2(\sinh T\pm T)>0$.  Applying
Theorem~\ref{thm:newdMovingConstraint} to
$L^2(I_{T'})=L^2(I_T)\oplus L^2(I_{T'}\setminus I_T)$ gives an explicit
orthogonal update of the primitive space $\ker B_T$.  At the arithmetic
thresholds
\[
 T_N=\tfrac12\log N,
 \qquad
 \delta_N=T_{N+1}-T_N=\tfrac12\log\frac{N+1}{N},
\]
the two nonzero eigenvalues of $C^*C$ are
\begin{equation}\label{eq:newdTateReturnEigenvalues}
 \kappa_{N,\pm}=
 \frac{\sinh T_{N+1}-\sinh T_N\ \pm\delta_N}
      {\sinh T_N\ \pm T_N}.
\end{equation}
Therefore
\begin{equation}\label{eq:newdTateReturnDecay}
 \|C\|^2=\frac1{2N}+o(N^{-1}),
 \qquad
 \|CM^{-1/2}\|=O(N^{-1/2}),
 \qquad
 \|M^{-1/2}-I\|=O(N^{-1}).
\end{equation}
Thus the motion of the two Tate constraints is an explicit conservative
rank-two terminal; it creates no unidentified projection term.

We next identify the finite incidence column with its physical
old/new boundary action.  Put $I_N=(-T_N,T_N)$ and
\[
 \mathcal E_N=
 (-T_{N+1},-T_N)\cup(T_N,T_{N+1}).
\]
Let $U_n$ be translation by $\log n$ on the zero-extension line and set
\begin{equation}\label{eq:newdPhysicalIncidence}
 \mathcal B_N^{\rm ar}
 =-P_{I_N}\sum_{2\leq n\leq N}
       \frac{\Lambda(n)}{\sqrt n}(U_n+U_n^*)
       \big|_{L^2(\mathcal E_N)}.
\end{equation}
After separating the two newborn wings, the translated intervals give an
isometry
\[
 \mathfrak S_N:
 \ell^2\{2,\ldots,N\}\widehat\otimes L^2(\mathcal E_N)
 \longrightarrow L^2(I_N),
\]
where the $n$th coordinate is the unique orientation entering $I_N$.

\begin{proposition}[Physical incidence fibre]
\label{prop:newdPhysicalIncidence}
If
\[
 \mathbf c_N=\sum_{n=2}^N\frac{\Lambda(n)}{\sqrt n}e_n,
\]
then, for every newborn profile $e$,
\begin{equation}\label{eq:newdPhysicalIncidenceColumn}
 \mathcal B_N^{\rm ar}e=-\mathfrak S_N(\mathbf c_N\otimes e),
\end{equation}
and
\begin{equation}\label{eq:newdPhysicalIncidenceGram}
 (\mathcal B_N^{\rm ar})^*\mathcal B_N^{\rm ar}
 =\left(\sum_{n=2}^N\frac{\Lambda(n)^2}{n}\right)I.
\end{equation}
At word depth $k$, the corresponding label column is
\begin{equation}\label{eq:newdWordColumn}
 (\log N)^{-k}\mathscr C_N^k\delta_1
 =\left\{\frac{\Lambda_k(n)}
                {\sqrt n(\log N)^k}\right\}_{n\leq N}.
\end{equation}
\end{proposition}

\begin{proof}
For a vector on the right newborn wing only $U_n$ enters the old interval;
for a vector on the left wing only $U_n^*$ enters it.  These translated
copies are pairwise disjoint, proving
\eqref{eq:newdPhysicalIncidenceColumn}; isometry gives
\eqref{eq:newdPhysicalIncidenceGram}.  Multiplication of an integer label by
$m$ agrees, whenever the result remains in the threshold fibre, with
physical translation by $\log m$.  Iterating and using Dirichlet convolution
gives \eqref{eq:newdWordColumn}.
\end{proof}

\subsubsection{Euler incidence and rational leakage}

On the right-oriented threshold fibre label the interval
$T_N-\log m+(0,\delta_N)$ by $m\leq N$.  Translation by $\log n$ sends
$m$ to $mn$, while translation in the opposite direction sends it to
$m/n$.  The latter remains in the integer fibre exactly when $n\mid m$.

\begin{proposition}[Divisibility and rational-leakage decomposition]
\label{prop:newdLeakageSplit}
The finite Euler boundary action is the orthogonal sum of its action on the
integer divisibility fibre and an action supported on nonintegral rational
labels.  An integer-labelled cell never overlaps a noninteger-labelled cell
in positive measure.  Two rational cells labelled by $m/n$ and $m'/n'$ can
overlap only when
\begin{equation}\label{eq:newdOverlapCondition}
 \left|\log\frac{m/n}{m'/n'}\right|<\delta_N.
\end{equation}
\end{proposition}

\begin{proof}
The assertions about labels follow directly from translation of the cell
positions.  If $m>kn$, then
\[
 \frac{m}{kn}\geq\frac N{N-1}>\frac{N+1}{N}=e^{2\delta_N};
\]
if $m<kn$, then $kn/m\geq(N+1)/N$, with equality producing at most an
endpoint contact.  Hence integer and noninteger cells are orthogonal.
Two intervals of length $\delta_N$ overlap exactly under
\eqref{eq:newdOverlapCondition}.
\end{proof}

Write $\Lambda_k=\Lambda^{*k}$.  For coprime positive integers $a,b$ with
$b>1$ define
\begin{equation}\label{eq:newdLeakageCoefficient}
 \ell_{N,k}(a,b)=
 \frac1{\sqrt{ab}(\log N)^k}
 \sum_{d\leq N/\max(a,b)}
       \frac{\Lambda(bd)\Lambda_k(ad)}d.
\end{equation}
Let
\[
 \mathscr R_N^{\rm leak}
 =\{a/b:(a,b)=1,\ b>1,\ \max(a,b)\leq N\},
\]
extend $e\in E_N:=L^2(0,\delta_N)$ by zero, and define
\begin{equation}\label{eq:newdLeakageOperator}
 L_{N,k}e=\sum_{a/b\in\mathscr R_N^{\rm leak}}
   \ell_{N,k}(a,b)\,\tau_{\log(a/b)}e,
 \qquad
 \mathbf L_Ne=\bigoplus_{k\geq1}L_{N,k}e.
\end{equation}
Only finitely many depths occur for a fixed $N$.

\begin{proposition}[Exact aggregation on reduced rational labels]
\label{prop:newdRationalAggregation}
The coefficient of the oppositely oriented Euler output on the reduced
rational cell $a/b$ is the sum in
\eqref{eq:newdLeakageCoefficient}.  If
$\kappa_e(u)=\langle\tau_ue,e\rangle$, then
\begin{align}
 \|\mathbf L_Ne\|^2
 ={}&\sum_{k\geq1}
 \sum_{a/b,c/d\in\mathscr R_N^{\rm leak}}
 \ell_{N,k}(a,b)\overline{\ell_{N,k}(c,d)}
 \kappa_e\!\left(\log\frac{ad}{bc}\right).
 \label{eq:newdLeakageGram}
\end{align}
Thus the Gram retains all coherent collisions between equal or nearby
reduced fractions.
\end{proposition}

\begin{proof}
The equality $m/n=a/b$ with $(a,b)=1$ is equivalent to
$m=ad$, $n=bd$ for a unique positive integer $d$.  Substitution of the
Euler weight and of the $k$th word column gives
\eqref{eq:newdLeakageCoefficient}.  Expanding the squared norm in
\eqref{eq:newdLeakageOperator} gives
\eqref{eq:newdLeakageGram}, since the inner product of the cells labelled by
$a/b$ and $c/d$ is $\kappa_e(\log(ad/bc))$.
\end{proof}

\subsubsection{Positive Gamma storage and the two exterior Tate states}

For $a>0$ and $f\in C_c^\infty(\R)$, let $z_{a,f}$ be the causal state
\begin{equation}\label{eq:newdGammaState}
 z_{a,f}'=-az_{a,f}+\sqrt{2a}\,f,
 \qquad z_{a,f}(-\infty)=0,
\end{equation}
and set $g_{a,f}=\sqrt{2a}\,z_{a,f}-f$.

\begin{lemma}[Lossless Gamma state]
\label{lem:newdGammaState}
The state satisfies
\begin{align}
 \frac d{dt}|z_{a,f}|^2&=|f|^2-|g_{a,f}|^2,
 \label{eq:newdGammaConservation}\\
 \frac1{2a}\|f-g_{a,f}\|_2^2
 &=\frac1{a^2}\int_\R|z_{a,f}'(t)|^2\,dt
 =\int_\R\frac{2\tau^2}{a(a^2+\tau^2)}
       |\widehat f(\tau)|^2\frac{d\tau}{2\pi}.
 \label{eq:newdGammaEnergy}
\end{align}
\end{lemma}

\begin{proof}
Expanding $|g_{a,f}|^2$ and using
\eqref{eq:newdGammaState} proves
\eqref{eq:newdGammaConservation}.  Solving
\eqref{eq:newdGammaState} for $f$ gives
$f-g_{a,f}=\sqrt{2/a}\,z_{a,f}'$.  Its Fourier multiplier has squared
modulus $4\tau^2/(a^2+\tau^2)$, proving
\eqref{eq:newdGammaEnergy}.
\end{proof}

Put $a_j=2j+\tfrac12$ and define the unreset Gamma energy
\begin{equation}\label{eq:newdFullGammaEnergy}
 \mathcal G_\Gamma[f]
 =\sum_{j\geq0}\frac1{a_j^2}\int_\R|z_{a_j,f}'(t)|^2\,dt
 =\int_\R g_\Gamma(\tau)|\widehat f(\tau)|^2
      \frac{d\tau}{2\pi},
\end{equation}
where
\begin{equation}\label{eq:newdGammaMultiplier}
 g_\Gamma(\tau)=
 \sum_{j\geq0}\frac{2\tau^2}{a_j(a_j^2+\tau^2)}.
\end{equation}
The causal equations are solved once on the real line; states are not reset
at internal rational-cell endpoints.

For a compactly supported $f$ define
\begin{equation}\label{eq:newdTateMoments}
 M_+(f)=\int_\R e^{t/2}f(t)\,dt,
 \qquad
 M_-(f)=\int_\R e^{-t/2}f(t)\,dt.
\end{equation}
If $I=(L,R)$ contains the support of $f$, the Gram of the two Tate vectors is
\begin{equation}\label{eq:newdExteriorTateGram}
 G_I=
 \begin{pmatrix}
 e^R-e^L&R-L\\
 R-L&e^{-L}-e^{-R}
 \end{pmatrix}>0,
\end{equation}
and the squared norm of the Tate component is
\begin{equation}\label{eq:newdTateTerminalForm}
 \mathcal T_I[f]=
 \begin{pmatrix}\overline{M_+(f)}&\overline{M_-(f)}\end{pmatrix}
 G_I^{-1}
 \binom{M_+(f)}{M_-(f)}.
\end{equation}
The moments in \eqref{eq:newdTateMoments} are exactly the exterior states of
\eqref{eq:newdGammaState} at $a=1/2$: for
$\operatorname{supp}f\subset[L,R]$,
\[
 z_{1/2,f}(t)=e^{-t/2}M_+(f)\quad(t\geq R),
\]
up to the normalization $\sqrt{2a}=1$, and the anticausal state gives
$e^{t/2}M_-(f)$ on $t\leq L$.

Take one interval containing every cell in
\eqref{eq:newdLeakageOperator}, for instance
\begin{equation}\label{eq:newdGlobalLeakageInterval}
 I_N=\left(-\log N,\ \log\frac N2+\delta_N\right).
\end{equation}
The source-defined Gamma--Tate form on the physical leakage column is
\begin{equation}\label{eq:newdGammaTateForm}
 \mathcal E_{\Gamma T,N}(e)=
 \sum_{k\geq1}
 \bigl\{\mathcal G_\Gamma[L_{N,k}e]
       +\mathcal T_{I_N}[L_{N,k}e]\bigr\}.
\end{equation}
Every term in this form is nonnegative and is determined by prime-power
weights, the Gamma parameters and the two Tate moments.  In particular, it
does not use the sign of the primitive Weil form.

\paragraph{The first nontrivial arithmetic cell.}
The threshold $N=3$ gives
\begin{equation}\label{eq:newdFirstCellLength}
 T_3=\tfrac12\log3,
 \qquad T_4=\log2,
 \qquad L_3=\delta_3=\tfrac12\log\frac43.
\end{equation}
At word depth one the only nonzero nonintegral coefficient is
\begin{equation}\label{eq:newdFirstLeakageCoefficient}
 \ell_{3,1}(2,3)
 =\frac{\Lambda(3)\Lambda(2)}{\sqrt6\log3}
 =\frac{\log2}{\sqrt6}.
\end{equation}
There is no deeper word, because $\Lambda_k(m)=0$ for $m\leq3$ and
$k\geq2$.  Consequently the complete rational leakage in this threshold
cell is
\begin{equation}\label{eq:newdFirstLeakageOutput}
 \mathbf L_3e=f_e
 :=\frac{\log2}{\sqrt6}\,
   \tau_{\log(2/3)}e,
 \qquad
 \|f_e\|_2^2=\frac{(\log2)^2}{6}\|e\|_2^2.
\end{equation}

\begin{theorem}[Complete Gamma--Tate payment on one arithmetic cell]
\label{thm:newdFirstCellPayment}
For every $e\in L^2(0,L_3)$, let $f_e$ be given by
\eqref{eq:newdFirstLeakageOutput}.  Then
\begin{equation}\label{eq:newdFirstCellPayment}
 \|\mathbf L_3e\|_2^2
 \leq \mathcal G_\Gamma[f_e]+\mathcal T_{J_3}[f_e],
 \qquad
 J_3=\log(2/3)+(0,L_3).
\end{equation}
More precisely, if $\mathcal G_{0}[f]$ denotes the first Gamma-channel
energy, corresponding to $a_0=1/2$, then
\begin{equation}\label{eq:newdFirstCellStrictPayment}
 \|f_e\|_2^2
 \leq \alpha_3\mathcal G_0[f_e]
       +\eta_3\mathcal T_{J_3}[f_e],
\end{equation}
where
\begin{equation}\label{eq:newdFirstCellConstants}
 \alpha_3=\frac14+\frac{L_3^2}{32}<0.251,
 \qquad
 \eta_3=\frac{1-e^{-L_3}}2
       =\frac{2-\sqrt3}{4}<0.067.
\end{equation}
Thus the Gamma interior energy and the two exterior Tate states pay the
whole leakage simultaneously, with a strict positive remainder:
\begin{align}
 &\mathcal G_\Gamma[f_e]+\mathcal T_{J_3}[f_e]
       -\|\mathbf L_3e\|_2^2                              \notag\\
 &\qquad\geq
 (1-\alpha_3)\mathcal G_0[f_e]
 +(1-\eta_3)\mathcal T_{J_3}[f_e]\geq0.
 \label{eq:newdFirstCellRemainder}
\end{align}
\end{theorem}

\begin{proof}
Translation does not change any of the following estimates, so write
$J_3=(0,L_3)$.  Let $z$ be the causal state of the first Gamma channel,
\[
 z'=-\tfrac12z+f_e,
 \qquad z(0)=0.
\]
The initial value vanishes because $f_e$ is supported on this single cell.
Since $f_e=z'+\tfrac12z$, integration of the cross term gives the exact
identity
\begin{equation}\label{eq:newdFirstCellStateIdentity}
 \|f_e\|_2^2
 =\|z'\|_2^2+\frac14\|z\|_2^2+\frac12|z(L_3)|^2.
\end{equation}
The anchored Cauchy--Schwarz inequality yields
\[
 |z(t)|^2
 \leq t\int_0^t|z'(s)|^2\,ds
 \leq t\|z'\|_2^2,
\]
and hence
\begin{equation}\label{eq:newdFirstCellPoincare}
 \|z\|_2^2\leq\frac{L_3^2}{2}\|z'\|_2^2.
\end{equation}
By \eqref{eq:newdGammaEnergy} at $a=1/2$,
\begin{equation}\label{eq:newdFirstCellGammaState}
 \mathcal G_0[f_e]=4\|z'\|_2^2.
\end{equation}

It remains to pay the endpoint.  The exterior-state formula gives
\[
 z(L_3)=e^{-L_3/2}M_+(f_e).
\]
The Tate terminal is the squared norm of the orthogonal projection onto
$\operatorname{span}\{e^{t/2},e^{-t/2}\}$.  It therefore dominates the
projection onto the first of these one-dimensional subspaces:
\begin{equation}\label{eq:newdFirstCellTateProjection}
 \mathcal T_{J_3}[f_e]
 \geq\frac{|M_+(f_e)|^2}{\|e^{t/2}\|_{L^2(J_3)}^2}
 =\frac{|M_+(f_e)|^2}{e^{L_3}-1}.
\end{equation}
Consequently
\begin{equation}\label{eq:newdFirstCellEndpoint}
 \frac12|z(L_3)|^2
 \leq\frac{1-e^{-L_3}}2\,\mathcal T_{J_3}[f_e]
 =\eta_3\mathcal T_{J_3}[f_e].
\end{equation}
Substitution of \eqref{eq:newdFirstCellPoincare},
\eqref{eq:newdFirstCellGammaState} and
\eqref{eq:newdFirstCellEndpoint} into
\eqref{eq:newdFirstCellStateIdentity} proves
\eqref{eq:newdFirstCellStrictPayment}.  Since the full Gamma energy contains
$\mathcal G_0$ as its first nonnegative summand and both constants in
\eqref{eq:newdFirstCellConstants} are smaller than one,
\eqref{eq:newdFirstCellPayment} and
\eqref{eq:newdFirstCellRemainder} follow.
\end{proof}

This theorem proves the complete one-cell assertion without a sign
assumption and without resetting a state inside the cell.  Its causal state
starts at zero only because, at $N=3$, the whole rational leakage consists
of the single support interval $J_3$.  For a connected cluster at a later
threshold, the outgoing state of one rational cell is the incoming state of
the next; Theorem~\ref{thm:newdFirstCellPayment} cannot be iterated by
setting those intermediate states to zero.

The same argument in fact applies to a whole connected cluster, provided it
is applied once on the exterior hull rather than separately on its
constituent cells.

\begin{lemma}[Gamma--Tate payment on a finite support hull]
\label{lem:newdClusterHull}
Let $f\in L^2(\R)$ be supported in an interval $J=(r,r+D)$.  If
$D<\sqrt{24}$, then
\begin{equation}\label{eq:newdClusterHullPayment}
 \|f\|_2^2
 \leq \alpha(D)\mathcal G_0[f]
       +\eta(D)\mathcal T_J[f]
 \leq \mathcal G_\Gamma[f]+\mathcal T_J[f],
\end{equation}
where
\begin{equation}\label{eq:newdClusterHullConstants}
 \alpha(D)=\frac14+\frac{D^2}{32}<1,
 \qquad
 \eta(D)=\frac{1-e^{-D}}2<1.
\end{equation}
No disjointness or orthogonality assumption is made inside $J$.
\end{lemma}

\begin{proof}
Translate $J$ to $(0,D)$ and solve the first causal Gamma equation once on
the complete hull:
\[
 z'=-\tfrac12z+f,
 \qquad z(0)=0.
\]
There is no condition on $z$ at an internal subdivision.  The proof of
Theorem~\ref{thm:newdFirstCellPayment}, with $D$ in place of $L_3$, gives
\begin{align*}
 \|f\|_2^2
 &=\|z'\|_2^2+\tfrac14\|z\|_2^2+\tfrac12|z(D)|^2,\\
 \|z\|_2^2&\leq\tfrac{D^2}{2}\|z'\|_2^2,
 \qquad
 \mathcal G_0[f]=4\|z'\|_2^2,\\
 \tfrac12|z(D)|^2
 &\leq\tfrac{1-e^{-D}}2\mathcal T_J[f].
\end{align*}
Combining these identities proves the first inequality in
\eqref{eq:newdClusterHullPayment}.  The restriction $D<\sqrt{24}$ is
exactly $\alpha(D)<1$; the second inequality then follows because the full
Gamma form contains $\mathcal G_0$ as a nonnegative summand.
\end{proof}

\paragraph{A three-cell connected cluster.}
Take $N=10$, word depth $k=1$, and
\begin{equation}\label{eq:newdThreeCellLabels}
 q_1=\frac34,
 \qquad q_2=\frac79,
 \qquad q_3=\frac45,
 \qquad
 \delta_{10}=\frac12\log\frac{11}{10}.
\end{equation}
The three coefficients obtained from
\eqref{eq:newdLeakageCoefficient} are
\begin{equation}\label{eq:newdThreeCellCoefficients}
 c_1=\frac{\log2\log3}{\sqrt{12}\log10},
 \qquad
 c_2=\frac{\log3\log7}{\sqrt{63}\log10},
 \qquad
 c_3=\frac{\log5\log2}{\sqrt{20}\log10}.
\end{equation}
They are all nonzero.  Moreover
\begin{equation}\label{eq:newdThreeCellConnections}
 \log\frac{q_2}{q_1}=\log\frac{28}{27}<\delta_{10},
 \qquad
 \log\frac{q_3}{q_2}=\log\frac{36}{35}<\delta_{10},
\end{equation}
so the translated cells overlap in a connected chain.  For
$e\in L^2(0,\delta_{10})$ define their coherent physical output
\begin{equation}\label{eq:newdThreeCellOutput}
 f_{10,e}=\sum_{j=1}^3c_j\tau_{\log q_j}e.
\end{equation}
Its support is contained in the single exterior hull
\begin{equation}\label{eq:newdThreeCellHull}
 J_{10}=\left(\log\frac34,
              \log\frac45+\delta_{10}\right),
 \qquad
 D_{10}=|J_{10}|=\log\frac{16}{15}
                 +\frac12\log\frac{11}{10}.
\end{equation}

\begin{theorem}[Complete payment of a connected three-cell cluster]
\label{thm:newdThreeCellCluster}
For every $e\in L^2(0,\delta_{10})$, including profiles for which the
adjacent translated copies interfere coherently,
\begin{equation}\label{eq:newdThreeCellPayment}
 \|f_{10,e}\|_2^2
 \leq\mathcal G_\Gamma[f_{10,e}]
      +\mathcal T_{J_{10}}[f_{10,e}].
\end{equation}
More precisely,
\begin{equation}\label{eq:newdThreeCellStrictPayment}
 \|f_{10,e}\|_2^2
 \leq\alpha_{10}\mathcal G_0[f_{10,e}]
      +\eta_{10}\mathcal T_{J_{10}}[f_{10,e}],
\end{equation}
where
\begin{align}
 \alpha_{10}
 &=\frac14+\frac1{32}
   \left(\log\frac{16}{15}
        +\frac12\log\frac{11}{10}\right)^2
   <0.251,\notag\\
 \eta_{10}
 &=\frac12\left[1-
   \exp\!\left(-\log\frac{16}{15}
               -\frac12\log\frac{11}{10}\right)\right]
   <0.054.                                                   
 \label{eq:newdThreeCellConstants}
\end{align}
Consequently
\begin{align}
 &\mathcal G_\Gamma[f_{10,e}]
  +\mathcal T_{J_{10}}[f_{10,e}]-\|f_{10,e}\|_2^2
 \notag\\
 &\qquad\geq
 (1-\alpha_{10})\mathcal G_0[f_{10,e}]
 +(1-\eta_{10})\mathcal T_{J_{10}}[f_{10,e}]\geq0.
 \label{eq:newdThreeCellRemainder}
\end{align}
\end{theorem}

\begin{proof}
The coefficient formulas follow because, for each pair in
\eqref{eq:newdThreeCellLabels}, the only contributing divisor in
\eqref{eq:newdLeakageCoefficient} is $d=1$.  Equations
\eqref{eq:newdThreeCellConnections} prove connectedness, and the extreme
endpoints give \eqref{eq:newdThreeCellHull}.  Numerically,
\[
 D_{10}=0.1121936110\ldots<\sqrt{24},
 \quad
 \alpha_{10}=0.2503933564\ldots,
 \quad
 \eta_{10}=0.0530644113\ldots.
\]
Apply Lemma~\ref{lem:newdClusterHull} directly to the coherent sum
$f_{10,e}$, before expanding its squared norm.  Therefore every diagonal
term and every interference term
\[
 c_i\overline{c_j}
 \langle\tau_{\log q_i}e,\tau_{\log q_j}e\rangle
\]
is included automatically.  The Gamma state is solved once from the left
endpoint of $J_{10}$ to its right endpoint and is never reset at either
internal cell boundary.  This proves
\eqref{eq:newdThreeCellStrictPayment}; the remaining statements follow as
in Lemma~\ref{lem:newdClusterHull}.
\end{proof}

Thus the earlier warning is not an obstruction to finite clusters: it only
forbids proving a cluster estimate by summing independently reset one-cell
estimates.  The exterior-hull argument proves every connected cluster with
hull length $D<\sqrt{24}$.  The following arithmetic separation lemma shows
that this restriction is automatic for every connected rational-leakage
cluster.  The point is that the cells respect a fixed family of reciprocal
half-integer bands.

\begin{theorem}[Reciprocal half-integer band theorem]
\label{thm:newdHalfIntegerBands}
Fix $N\geq2$ and a word depth $k$.  Form the intervals
\[
 J_q=\log q+(0,\delta_N),
 \qquad
 q=a/b\in\mathscr R_N^{\rm leak},
 \qquad \ell_{N,k}(a,b)>0.
\]
Every connected component of $\bigcup_qJ_q$ is contained in one of the
bands
\begin{equation}\label{eq:newdArithmeticBands}
 \left(\log\frac m2,\log\frac{m+1}{2}\right)
 \quad(m\geq2),
 \qquad\text{or}\qquad
 \left(\log\frac2{m+1},\log\frac2m\right)
 \quad(m\geq2),
\end{equation}
up to the possible inclusion of one boundary cell.  In particular, the
exterior hull of every connected component has length
\begin{equation}\label{eq:newdUniformClusterDiameter}
 D_{N,k}^{\rm comp}
 \leq \log\frac32+\delta_N
 \leq\frac32\log\frac32
 <\sqrt{24}.
\end{equation}
Consequently every connected component $C$ and every profile
$e\in L^2(0,\delta_N)$ satisfy, for
\[
 f_{C,e}=\sum_{q:\,J_q\subset C}
          \ell_{N,k}(q)\tau_{\log q}e,
\]
the complete estimate
\begin{equation}\label{eq:newdEveryConnectedClusterPaid}
 \|f_{C,e}\|_2^2
 \leq\mathcal G_\Gamma[f_{C,e}]
      +\mathcal T_{\operatorname{hull}(C)}[f_{C,e}].
\end{equation}
This assertion is uniform in $N$, $k$, the number of cells in $C$, and
all coherent interference among them.
\end{theorem}

\begin{proof}
We first prove that a rational cell cannot cross the indicated barriers.
Let $q=a/b>1$ be reduced, with $a,b\leq N$, and let $r=m/2>q$ be a
half-integer.  Since $mb-2a$ is a positive integer,
\[
 \log\frac rq
 =\log\frac{mb}{2a}
 \geq\log\left(1+\frac1{2a}\right)
 \geq\log\left(1+\frac1{2N}\right).
\]
But
\begin{equation}\label{eq:newdBarrierMargin}
 \log\left(1+\frac1{2N}\right)
 >\frac12\log\left(1+\frac1N\right)=\delta_N,
\end{equation}
because $(1+1/(2N))^2>1+1/N$.  Hence the right endpoint of $J_q$ lies
strictly to the left of $\log r$.

For $0<q=a/b<r=2/m$, the positive integer $2b-ma$ gives
\[
 \log\frac rq
 =\log\frac{2b}{ma}
 \geq\log\left(1+\frac1{ma}\right).
\]
Here $ma=2b-(2b-ma)\leq2N-1$, and therefore the last quantity is at least
$\log(1+1/(2N-1))$, which is strictly larger than the left side of
\eqref{eq:newdBarrierMargin}, hence larger than $\delta_N$.  Thus no cell
below a reciprocal barrier $2/m$ crosses it either.  Cells starting on a
barrier belong only to the band on its right, so they do not join the two
adjacent open bands.  Notice that this proof uses only $a,b\leq N$; it is
therefore valid at every word depth and does not depend on which of the
coefficients happen to vanish.

The consecutive ratios of the band endpoints are
\[
 \frac{(m+1)/2}{m/2}=1+\frac1m,
 \qquad
 \frac{2/m}{2/(m+1)}=1+\frac1m.
\]
Their logarithms are at most $\log(3/2)$ for $m\geq2$.  Adding one cell
length gives the first inequality in
\eqref{eq:newdUniformClusterDiameter}; the second follows from
$\delta_N\leq\delta_2=\tfrac12\log(3/2)$.  Lemma
\ref{lem:newdClusterHull} now applies once to the complete coherent output
on each connected component and proves
\eqref{eq:newdEveryConnectedClusterPaid}.
\end{proof}

\begin{corollary}[No low-frequency fugitive inside one arithmetic band]
\label{cor:newdBandCoherence}
For a connected component $C$ define its positive exponential polynomial
\[
 P_C(\tau)=\sum_{q:\,J_q\subset C}\ell_{N,k}(q)q^{-i\tau}.
\]
For every $0\leq\omega<\pi/\log(3/2)$ and $|\tau|\leq\omega$ one has
\begin{equation}\label{eq:newdBandCoherence}
 c(\omega)P_C(0)\leq |P_C(\tau)|\leq P_C(0),
 \qquad
 c(\omega)=
 \cos\!\left(\frac{\omega}{2}\log\frac32\right)>0.
\end{equation}
Thus the translates belonging to a single arithmetic band cannot, at any
$N$ or word depth, form a low-frequency fugitive profile by destructive
interference.
\end{corollary}

\begin{proof}
Let $x_q=\log q$ and let $x_C$ be the midpoint of the smallest interval
containing all $x_q$ in $C$.  The band theorem gives
$|x_q-x_C|\leq\tfrac12\log(3/2)$.  Positivity of the coefficients yields
\begin{align*}
 |P_C(\tau)|
 &\geq \Re\!\left(e^{i\tau x_C}P_C(\tau)\right)\\
 &=\sum_{q\in C}\ell_{N,k}(q)
       \cos\bigl(\tau(x_q-x_C)\bigr)\\
 &\geq
 \cos\!\left(\frac{\omega}{2}\log\frac32\right)
 \sum_{q\in C}\ell_{N,k}(q).
\end{align*}
The cosine is positive by the assumed range of $\omega$.  The upper bound
is the triangle inequality.
\end{proof}

For orientation, direct enumeration of the exact nonzero coefficients at
depth one gives the following largest connected hulls.  These figures are
not used in the proof; they display the approach to the sharp first-band
width $\log(3/2)$.
\[
\begin{array}{c|c|c|c}
N&\text{leftmost label}&\text{rightmost label}&
 D_{N,1}^{\rm largest}\\ \hline
10&6/5&4/3&0.1530156\ldots\\
50&33/49&48/49&0.3845948\ldots\\
200&133/199&198/199&0.4004117\ldots\\
1000&665/997&996/997&0.4044600\ldots
\end{array}
\]

The theorem removes the possibility of an arbitrarily long connected
cluster, and Corollary~\ref{cor:newdBandCoherence} removes destructive
low-frequency interference inside each such cluster.  It does not by
itself permit one to add
\eqref{eq:newdEveryConnectedClusterPaid} over disconnected bands: the
unreset Gamma form and the two global Tate moments contain cross-band
terms.  The remaining question is therefore a recombination theorem for
the separated arithmetic bands---equivalently, exclusion of cancellation
between distinct band polynomials---not a further one-cell or
connected-cluster estimate.

\paragraph{Audit of causal recombination across the gaps.}
There is a tempting but incorrect argument which treats decay of the causal
state in a gap as a negative Schur remainder.  The exact sign is the
opposite.  To record this point unambiguously, write
\[
 A_m=z(\inf C_m),\qquad B_m=z(\sup C_m),\qquad
 A_{m+1}=e^{-\Delta_m/2}B_m,qquad A_1=0.
\]
Summing the exact cluster identities gives
\begin{align}
 \frac12\sum_{m=1}^M(|B_m|^2-|A_m|^2)
 &=\frac12|B_M|^2
   +\frac12\sum_{m=1}^{M-1}
      (1-e^{-\Delta_m})|B_m|^2.             
 \label{eq:newdGapAuditSign}
\end{align}
Thus the gap term occurs with a plus sign, not a minus sign.  Equivalently,
on a gap the equation $z'=-z/2$ gives
\[
 4\int_{\rm gap}|z'|^2
 =\int_{\rm gap}|z|^2
 =(1-e^{-\Delta_m})|B_m|^2,
\]
and the two positive bulk terms are cancelled exactly by the decrease of
the endpoint storage when the identity is written on the complete hull;
they do not produce an additional negative remainder.  Indeed the global
identity remains
\begin{equation}\label{eq:newdGlobalStateAuditIdentity}
 \|f\|_2^2
 =\frac14\mathcal G_0[f]
  +\frac14\|z\|_2^2+rac12|B_M|^2.
\end{equation}

There is a second independent obstruction to that argument.  The anchored
estimate
\[
 \|z\|_{L^2(C_m)}^2
 \leq\frac{|C_m|^2}{2}\|z'\|_{L^2(C_m)}^2
\]
used in Lemma~\ref{lem:newdClusterHull} requires $A_m=0$ and cannot be
reused for the restriction of the global state to later clusters.  For
example, the homogeneous state $z(t)=Ae^{-(t-L)/2}$ on an interval of
length $D$ has
\[
 \|z'\|_2^2=\frac14\|z\|_2^2,
\]
so the displayed estimate would require $1\leq D^2/8$, which is false for
all the band lengths in \eqref{eq:newdUniformClusterDiameter}.  An
arbitrarily small nonzero source perturbation gives the same failure inside
a genuine later support component.  Therefore causal decay by itself does
not prove the required recombination; a valid proof must control the
incoming states or establish the global low-frequency observability
inequality below.

\paragraph{The single obstruction for arbitrary clusters.}
The preceding results permit an exact separation of the disconnected-band
problem.  Since every summand in \eqref{eq:newdGammaMultiplier} is strictly
increasing for $\tau>0$, while $g_\Gamma(0)=0$ and
$g_\Gamma(\tau)\to\infty$, there is a unique number
$\tau_\Gamma>0$ such that
\begin{equation}\label{eq:newdGammaCutoff}
 g_\Gamma(\tau_\Gamma)=1.
\end{equation}
Numerically, $\tau_\Gamma=0.2860886323\ldots$; the numerical value is not
used below.  Put $\Omega_\Gamma=(-\tau_\Gamma,\tau_\Gamma)$ and define the
positive and negative Gamma weights
\[
 w_-(\tau)=(1-g_\Gamma(\tau))_+,
 \qquad
 w_+(\tau)=(g_\Gamma(\tau)-1)_+.
\]
In fact $\tau_\Gamma<1$, because the $j=0$ summand alone gives
$g_\Gamma(1)>1$, and $1<\pi/\log(3/2)$.  Hence
Corollary~\ref{cor:newdBandCoherence} applies on the whole deficit window
$\Omega_\Gamma$: no individual reciprocal half-integer band has a
low-frequency zero or even an arbitrarily small normalized value there.
For the $k$th rational leakage output write
\begin{equation}\label{eq:newdLeakagePolynomial}
 P_{N,k}(\tau)=
 \sum_{a/b\in\mathscr R_N^{\rm leak}}
 \ell_{N,k}(a,b)e^{-i\tau\log(a/b)},
 \qquad
 \widehat{L_{N,k}e}(\tau)=\widehat e(\tau)P_{N,k}(\tau).
\end{equation}

\begin{proposition}[Exact low-frequency reduction]
\label{prop:newdLowFrequencyReduction}
The uniform Gamma--Tate estimate
\eqref{eq:newdUniformSourceEstimate} is equivalent to
\begin{align}
 &\sum_{k\geq1}\int_{\Omega_\Gamma}
    w_-(\tau)|\widehat e(\tau)P_{N,k}(\tau)|^2
       \frac{d\tau}{2\pi}                                \notag\\
 &\qquad\leq
 \sum_{k\geq1}\mathcal T_{I_N}[L_{N,k}e]
 +\sum_{k\geq1}\int_{\R\setminus\Omega_\Gamma}
    w_+(\tau)|\widehat e(\tau)P_{N,k}(\tau)|^2
       \frac{d\tau}{2\pi}.                               
 \label{eq:newdLowFrequencyObservability}
\end{align}
Moreover, because all coefficients
$\ell_{N,k}(a,b)$ are nonnegative,
\begin{equation}\label{eq:newdPositiveCoefficientEnvelope}
 |P_{N,k}(\tau)|^2
 \leq P_{N,k}(i/2)P_{N,k}(-i/2)
 \qquad(\tau\in\R).
\end{equation}
Thus every frequency outside the fixed compact band
$\Omega_\Gamma$ is already paid by Gamma.  The only remaining issue is the
two-moment observability of the arithmetic leakage range inside that band.
\end{proposition}

\begin{proof}
For $F_k=L_{N,k}e$, Plancherel and
\eqref{eq:newdFullGammaEnergy} give
\begin{align*}
 &\mathcal G_\Gamma[F_k]+\mathcal T_{I_N}[F_k]-\|F_k\|_2^2\\
 &\quad=\mathcal T_{I_N}[F_k]
 +\int_\R(g_\Gamma(\tau)-1)|\widehat F_k(\tau)|^2
          \frac{d\tau}{2\pi}.
\end{align*}
Splitting the integral at $\Omega_\Gamma$, substituting
\eqref{eq:newdLeakagePolynomial}, and summing over $k$ proves the
equivalence.  For the envelope, write $q=a/b$ and use Cauchy--Schwarz:
\begin{align*}
 |P_{N,k}(\tau)|
 &\leq\sum_q\ell_{N,k}(q)\\
 &=\sum_q
   \sqrt{\ell_{N,k}(q)q^{1/2}}
   \sqrt{\ell_{N,k}(q)q^{-1/2}}\\
 &\leq
 \left(\sum_q\ell_{N,k}(q)q^{1/2}\right)^{1/2}
 \left(\sum_q\ell_{N,k}(q)q^{-1/2}\right)^{1/2}.
\end{align*}
The two factors are $P_{N,k}(i/2)^{1/2}$ and
$P_{N,k}(-i/2)^{1/2}$, proving
\eqref{eq:newdPositiveCoefficientEnvelope}.
\end{proof}

\paragraph{The mean/oscillation Schur decomposition.}
The preceding band calculation suggests a sharper way to formulate the
remaining global problem.  Put
\[
 \mathfrak A_N[e]=\|\mathbf L_Ne\|_2^2,
 \qquad
 \mathfrak Q_N[e]=\mathcal E_{\Gamma T,N}(e),
 \qquad
 \mathfrak D_N=\mathfrak Q_N-\mathfrak A_N,
\]
and decompose the input cell exactly as
\begin{equation}\label{eq:newdMeanZeroSplit}
 E_N=\C u_N\oplus E_N^0,
 \qquad
 u_N=\delta_N^{-1/2}\mathbf1_{(0,\delta_N)},
 \qquad
 E_N^0=\left\{h:\int_0^{\delta_N}h=0\right\}.
\end{equation}
Let $d_N=\mathfrak D_N[u_N]$, let
$b_N(h)=\mathfrak D_N(u_N,h)$, and let $C_N$ be the restriction of
$\mathfrak D_N$ to $E_N^0$.

\begin{proposition}[Exact mean/oscillation recombination criterion]
\label{prop:newdMeanZeroSchur}
The unit Gamma--Tate estimate at threshold $N$ holds if and only if
\begin{equation}\label{eq:newdMeanZeroSchurCriterion}
 C_N\geq0,
 \qquad
 b_N\in\operatorname{Ran}C_N^{1/2},
 \qquad
 d_N-b_N^*C_N^\dagger b_N\geq0.
\end{equation}
Thus the remaining interference problem splits into a coercivity statement
for zero-mean cell profiles and one scalar Schur inequality for the constant
profile after all oscillatory profiles have been eliminated.
\end{proposition}

\begin{proof}
In the orthogonal decomposition \eqref{eq:newdMeanZeroSplit}, the defect
form has block matrix
\[
 \mathfrak D_N=
 \begin{pmatrix}d_N&b_N^*\\ b_N&C_N\end{pmatrix}.
\]
The generalized Schur-complement criterion for a closed Hermitian form is
exactly \eqref{eq:newdMeanZeroSchurCriterion}.  Positivity of
$\mathfrak D_N$ is the unit Gamma--Tate estimate by definition.
\end{proof}

\paragraph{Finite-dimensional reconnaissance (audit only).}
To identify which part of Proposition~\ref{prop:newdMeanZeroSchur} carries
the difficulty, the exact arithmetic coefficients were evaluated on the
space of twelve-bin piecewise-constant profiles on $(0,\delta_N)$.  The
table reports the smallest generalized quotient
$\mathfrak Q_N/\mathfrak A_N$, the same quotient after imposing zero mean,
and the proportion of the constant defect consumed by its Schur correction.
\[
\begin{array}{c|c|c|c}
N&\min \mathfrak Q_N/\mathfrak A_N&
\min_{E_N^0}\mathfrak Q_N/\mathfrak A_N&
b_N^*C_N^{-1}b_N/d_N\\ \hline
10 &5.1436&8.2823&0.0156\\
20 &4.9506&8.9386&0.0175\\
40 &4.8129&9.6305&0.0150\\
80 &4.5263&10.2950&0.0116\\
120&4.4563&10.6917&0.0118
\end{array}
\]
Mesh refinement at $N=80$ from six to sixteen bins changes the first
quotient from $4.5316$ to $4.5263$ and the zero-mean quotient from
$10.3691$ to $10.2807$.  These computations are neither interval-certified
nor a proof for arbitrary profiles.  Their consistent structural signal is
that the prospective minimizer lies in the almost-constant sector, while
the zero-mean sector has a large Gamma margin and the cross-sector Schur
correction is small.  Accordingly, the next analytic target is
\eqref{eq:newdMeanZeroSchurCriterion}, not an unstructured estimate over
all cell profiles at once.

The restriction to the arithmetic range in
\eqref{eq:newdLowFrequencyObservability} is essential.  Indeed, choose
$\varphi\in C_c^\infty(-1/4,1/4)$ with $\|\varphi\|_2=1$ and put
\begin{equation}\label{eq:newdPlateauFamily}
 f_D(t)=D^{-1/2}\varphi(t/D),
 \qquad I_D=(-D/2,D/2).
\end{equation}
Then $\|f_D\|_2=1$, whereas
\begin{equation}\label{eq:newdPlateauEscape}
 \mathcal G_\Gamma[f_D]\longrightarrow0,
 \qquad
 \mathcal T_{I_D}[f_D]\longrightarrow0
 \qquad(D\to\infty).
\end{equation}
For the first limit, substitute
$\widehat f_D(\tau)=D^{1/2}\widehat\varphi(D\tau)$ and use
$g_\Gamma(u/D)\to0$ together with the rapid decay of
$\widehat\varphi$.  For the second, the support of $f_D$ stays a distance
$D/4$ from both endpoints of $I_D$: its two exponential moments grow at
most like $D^{1/2}e^{D/8}$, while the corresponding diagonal Tate Grams
grow like $e^{D/2}$; the off-diagonal entry is only $D$.  Hence the Tate
projection tends to zero exponentially.

Consequently no support-only argument can prove the unit estimate for an
aggregate whose disconnected bands have an arbitrarily long exterior
hull.  The one remaining theorem is precisely that
outputs of the positive arithmetic polynomials
$P_{N,k}$ cannot form the escaping family
\eqref{eq:newdPlateauFamily}; equivalently, they must satisfy the compact
low-frequency observability inequality
\eqref{eq:newdLowFrequencyObservability} uniformly in $N$.

\subsubsection{The exact threshold Schur condition}

Let the two physical analysis factors at one threshold have old/new columns
\[
 X=(X_0,X_E),\qquad Y=(Y_0,Y_E),
\]
and put
\begin{equation}\label{eq:newdPhysicalBlocks}
 R_0=X_0^*X_0,
 \quad L_0=Y_0^*Y_0,
 \quad r=X_0^*X_E,
 \quad \ell=Y_0^*Y_E,
 \quad A_0=R_0-L_0.
\end{equation}
Assume the old primitive block $A_0$ is nonnegative and the reference range
condition $r\in\operatorname{Ran}R_0$ holds.  Define
\begin{align}
 H&=R_0^\dagger r,
 \label{eq:newdReferenceLift}\\
 S_E&=X_E^*X_E-r^*R_0^\dagger r,
 \label{eq:newdReferenceSchur}\\
 Z_E&=Y_E-Y_0H,
 \qquad
 b_E=L_0H-\ell.
 \label{eq:newdOppositeResidual}
\end{align}

\begin{theorem}[Exact threshold completion criterion]
\label{thm:newdThresholdCriterion}
The full old/new signed block $X^*X-Y^*Y$ is nonnegative if and only if
\begin{equation}\label{eq:newdRangeCondition}
 b_E\in\operatorname{Ran}A_0^{1/2}
\end{equation}
and
\begin{equation}\label{eq:newdSchurTarget}
 S_E-Z_E^*Z_E-b_E^*A_0^\dagger b_E\geq0.
\end{equation}
\end{theorem}

\begin{proof}
The equation $R_0H=r$ gives
$X_0^*(X_E-X_0H)=0$ and
\[
 (X_E-X_0H)^*(X_E-X_0H)=S_E.
\]
Hence the reference-harmonic column $v_E=(-H,I)^t$ satisfies
\[
 v_E^*(X^*X-Y^*Y)v_E=S_E-Z_E^*Z_E.
\]
Its old cross is
\[
 (r-\ell)-(R_0-L_0)H=L_0H-\ell=b_E.
\]
Congruence by the invertible triangular change from the coordinate newborn
column to $v_E$ leaves the old block equal to $A_0$ and changes the other
entries to $b_E$ and $S_E-Z_E^*Z_E$.  The generalized Schur-complement
criterion gives precisely \eqref{eq:newdRangeCondition} and
\eqref{eq:newdSchurTarget}.
\end{proof}

This theorem is the fixed global target.  The preceding construction
determines the divergence, the moving Tate cell, the complete positive Gamma
form and the exact rational leakage.  It does not determine the sign in
\eqref{eq:newdSchurTarget}.

The source construction isolates a concrete sufficient estimate.  Let
\[
 \mathscr D_N=
 \{e\in E_N:\mathcal E_{\Gamma T,N}(e)<\infty\}.
\]
It would suffice to prove, with the same unit constant for every arithmetic
threshold,
\begin{equation}\label{eq:newdUniformSourceEstimate}
 \|\mathbf L_Ne\|^2\leq\mathcal E_{\Gamma T,N}(e)
 \qquad(N\geq2,\ e\in\mathscr D_N).
\end{equation}
Indeed, Proposition~\ref{prop:newdLeakageSplit} separates the integer
incidence from the noninteger leakage.  The former is the finite
divisibility current of Theorem~\ref{thm:newdAmplitude}; the latter is
exactly $\mathbf L_Ne$ by
Proposition~\ref{prop:newdRationalAggregation}.  The right side of
\eqref{eq:newdUniformSourceEstimate} is the unreset positive Gamma storage
plus the conservative two-dimensional Tate terminal.  Thus the estimate
would make the joint boundary source contractive with unit normalization;
shorting a passive source network would give
\eqref{eq:newdRangeCondition}--\eqref{eq:newdSchurTarget} at the next
threshold.  Induction from the certified initial interval would then give
the primitive inequality for every $T$.

For clarity, \eqref{eq:newdUniformSourceEstimate} is a sufficient
source-level route, not a definition of the desired Schur complement and not
an additional conclusion of the identities already proved.  Expanded on the
test core $C_c^\infty(0,\delta_N)$, it is the explicit inequality
\begin{align}
 &\sum_{k\geq1}
 \sum_{a/b,c/d\in\mathscr R_N^{\rm leak}}
 \ell_{N,k}(a,b)\overline{\ell_{N,k}(c,d)}
 \kappa_e\!\left(\log\frac{ad}{bc}\right)
 \notag\\
 &\quad\leq
 \sum_{k\geq1}\int_\R g_\Gamma(\tau)
 \left|\widehat e(\tau)
   \sum_{a/b\in\mathscr R_N^{\rm leak}}
       \ell_{N,k}(a,b)e^{-i\tau\log(a/b)}\right|^2
 \frac{d\tau}{2\pi}
 \notag\\
 &\qquad+
 \sum_{k\geq1}
 \begin{pmatrix}
  \overline{M_+(L_{N,k}e)}&\overline{M_-(L_{N,k}e)}
 \end{pmatrix}
 G_{I_N}^{-1}
 \binom{M_+(L_{N,k}e)}{M_-(L_{N,k}e)}.
 \label{eq:newdExpandedSourceEstimate}
\end{align}
All quantities in \eqref{eq:newdExpandedSourceEstimate} have been defined
independently of the desired sign.  The unresolved issue is its uniform
validity, including coherent clusters of nearby rational labels; a cellwise
argument that resets the Gamma state does not prove it.

\begin{center}
\fbox{\parbox{0.91\linewidth}{
\textbf{Open end of row (d).}
For every threshold, prove the range condition
$b_E\in\operatorname{Ran}A_0^{1/2}$ and the uniform Schur inequality
\[
 S_E-Z_E^*Z_E-b_E^*A_0^\dagger b_E\geq0.
\]
Equivalently, construct from the explicit Gamma--Euler--Tate source data a
Hilbert-space factorization of the corresponding old/new block.  A concrete
sufficient route is the unit-constant estimate
\eqref{eq:newdUniformSourceEstimate}, or its expanded form
\eqref{eq:newdExpandedSourceEstimate}, for every $N$ and every admissible
cell profile.  This uniform estimate, the global primitive inequality, and
row~(d) are not proved in the present manuscript.
}}
\end{center}
